\section{Conclusion}\label{sec:conclusion}

In this study, we developed a computational pricing pipeline for American options under regime-coupled stochastic volatility, in which the IV-shape function was calibrated once against a historical option-ladder corpus and forward simulation then propagated path-conditional implied volatility without re-ingesting market IV along the simulated path. The pipeline linked a Jump Hidden Markov Model to a modified Heston variance process via a state-dependent mean-reversion target $\theta(t)$, and priced the resulting IV through American lattice pricers (a CRR binomial tree for the pricing pipeline, a Leisen-Reimer tree for the forward scenarios). The equilibrium initialization $v_0 = \theta(t\!=\!0)$ eliminated the initial IV surface as a free parameter at simulation time; the endogenous market mood signal was designed to produce correlated IV spikes during stress events without external data, although the reported experiments held the regime-state and mood multipliers inactive.

We calibrated the smile and term-structure shape function $\psi$ through a hierarchy of representations on a 31-ticker, fifteen-date, six-sector option ladder. The five-parameter parametric form served as the analytic baseline; replacing it with a globally shared neural surrogate and then with sector-specific neural surrogates reduced the overall fit error by 8\% and 18\% respectively (Table~\ref{tab:three_way}), with independent contributions from added model capacity and loosened sharing granularity. A temporal holdout on an eight-day subset of the same capture further showed that excluding scheduled earnings windows eliminated the test-time generalization gap of the sector NN, and that adding calendar-derived earnings-distance and same-sector peer-distance features to $\psi$ recovered roughly one fifth of that gap (Table~\ref{tab:earnings_holdout}), narrowing the holdout generalization gap from $+4.99\%$ to $+4.36\%$ IV; the residual $+4.36\%$ portion arose from earnings surprises whose direction and magnitude could not be inferred from the calendar alone. We then exercised the calibrated framework as a synthetic-data generator on a real GS 2026-05-29 \$890 put and \$970 call (both 31 DTE, 30-delta) pulled from the 04-28-2026 capture, forward-simulated 1{,}000 price paths under the JumpHMM and the per-ticker $\psi_{\mathrm{NN}}$, and tracked path-conditional implied volatility, premium decay, finite-difference Leisen-Reimer Greeks, and terminal short-premium P\&L from a single coherent simulation. The short put and short call kept their full premium on $74\%$ and $69\%$ of paths respectively, and the asymmetric tail risk of the short call (worst-case loss $1.57\times$ that of the short put despite a smaller premium received; Supplementary Table~\ref{tab:gs_scenario}, Fig.~\ref{fig:gs_pnl}) emerged naturally from the negative skew that the per-ticker $\psi_{\mathrm{NN}}$ encoded together with the leverage coupling that lifted $v_t$ on down-paths. A cross-ticker re-run on LLY (Fig.~\ref{fig:lly_pnl}, Supplementary Table~\ref{tab:lly_scenario}) reproduced the same qualitative asymmetry on a ticker with a different sector and a higher IV regime, confirming the result was structural rather than ticker-specific. The framework was implemented as an open-source Julia package and is available at \url{https://github.com/varnerlab/heston-implied-volatility-model}.
